% GENERATED FILE. Edit canonical lessons, then run npm run generate:books.
% canonical-source-hash: fnv1a64:7fb9f6f9119b137d
% canonical-lessons: JA-C89-kowasu, JA-C89-wakasu, JA-C89-mawasu, JA-C89-tsukamu, JA-C89-tatamu

\chapter{To break, To boil, To turn round, To grab, To fold}
\label{ch:ja-to-break-to-boil-to-turn-round-to-grab-to-fold}
\hlchaptermodality{89}

\begin{chapteropening}
\textbf{By the end of this chapter:} \emph{I can say to break, to boil, to turn round, to grab and to fold.}

Complete the last lesson of chapter 89: i can say to break, to boil, to turn round, to grab and to fold.
\end{chapteropening}

\section[\texorpdfstring{\ja{こわす} (kowasu)}{こわす (kowasu)}]{\ja{こわす} (kowasu) — to break}

\label{lesson:JA-C89-kowasu}



\begin{quote}
Before the new one: say the Japanese for to leave behind, then the Japanese for to knock down.
\end{quote}

\subsection*{You'll want to know: \ja{こわす}}
\textbf{\ja{こわす}} — \emph{kowasu} — \textquotedblleft{}to break\textquotedblright{}.

\begin{grammarlens}[title={What you've built}]
One more word for this chapter.
\end{grammarlens}

\subsection*{Your turn}
\begin{itemize}
\raggedright
  \item \emph{Say it:} \emph{kowasu}
  \item \emph{Say it:} \emph{kowasu}, once more
  \item \emph{Say it:} say \emph{kowasu}
  \item \emph{Recall:} say the Japanese for to leave behind, then the Japanese for to knock down, then say \emph{kowasu} again
\end{itemize}

\begin{tcolorbox}[breakable,colback=teal!4,colframe=teal!35!black,title={Before you move on}]
What does \ja{こわす} mean? (\textquotedblleft{}To break\textquotedblright{}.) Say it once more.
\end{tcolorbox}
\section[\texorpdfstring{\ja{わかす} (wakasu)}{わかす (wakasu)}]{\ja{わかす} (wakasu) — to boil (water)}

\label{lesson:JA-C89-wakasu}



\begin{quote}
Before the new one: say the Japanese for to knock down, then the Japanese for to break.
\end{quote}

\subsection*{You'll want to know: \ja{わかす}}
\textbf{\ja{わかす}} — \emph{wakasu} — \textquotedblleft{}to boil (water)\textquotedblright{}.

\begin{grammarlens}[title={What you've built}]
One more word for this chapter.
\end{grammarlens}

\subsection*{Your turn}
\begin{itemize}
\raggedright
  \item \emph{Say it:} \emph{wakasu}
  \item \emph{Say it:} \emph{wakasu}, once more
  \item \emph{Say it:} say \emph{wakasu}
  \item \emph{Recall:} say the Japanese for to knock down, then the Japanese for to break, then say \emph{wakasu} again
\end{itemize}

\begin{tcolorbox}[breakable,colback=teal!4,colframe=teal!35!black,title={Before you move on}]
What does \ja{わかす} mean? (\textquotedblleft{}To boil (water)\textquotedblright{}.) Say it once more.
\end{tcolorbox}
\section[\texorpdfstring{\ja{まわす} (mawasu)}{まわす (mawasu)}]{\ja{まわす} (mawasu) — to turn round}

\label{lesson:JA-C89-mawasu}



\begin{quote}
Before the new one: say the Japanese for to break, then the Japanese for to boil.
\end{quote}

\subsection*{You'll want to know: \ja{まわす}}
\textbf{\ja{まわす}} — \emph{mawasu} — \textquotedblleft{}to turn round\textquotedblright{}.

\begin{grammarlens}[title={What you've built}]
One more word for this chapter.
\end{grammarlens}

\subsection*{Your turn}
\begin{itemize}
\raggedright
  \item \emph{Say it:} \emph{mawasu}
  \item \emph{Say it:} \emph{mawasu}, once more
  \item \emph{Say it:} say \emph{mawasu}
  \item \emph{Recall:} say the Japanese for to break, then the Japanese for to boil, then say \emph{mawasu} again
\end{itemize}

\begin{tcolorbox}[breakable,colback=teal!4,colframe=teal!35!black,title={Before you move on}]
What does \ja{まわす} mean? (\textquotedblleft{}To turn round\textquotedblright{}.) Say it once more.
\end{tcolorbox}
\section[\texorpdfstring{\ja{つかむ} (tsukamu)}{つかむ (tsukamu)}]{\ja{つかむ} (tsukamu) — to grab}

\label{lesson:JA-C89-tsukamu}



\begin{quote}
Before the new one: say the Japanese for to boil, then the Japanese for to turn round.
\end{quote}

\subsection*{You'll want to know: \ja{つかむ}}
\textbf{\ja{つかむ}} — \emph{tsukamu} — \textquotedblleft{}to grab\textquotedblright{}.

\begin{grammarlens}[title={What you've built}]
One more word for this chapter.
\end{grammarlens}

\subsection*{Your turn}
\begin{itemize}
\raggedright
  \item \emph{Say it:} \emph{tsukamu}
  \item \emph{Say it:} \emph{tsukamu}, once more
  \item \emph{Say it:} say \emph{tsukamu}
  \item \emph{Recall:} say the Japanese for to boil, then the Japanese for to turn round, then say \emph{tsukamu} again
\end{itemize}

\begin{tcolorbox}[breakable,colback=teal!4,colframe=teal!35!black,title={Before you move on}]
What does \ja{つかむ} mean? (\textquotedblleft{}To grab\textquotedblright{}.) Say it once more.
\end{tcolorbox}
\section[\texorpdfstring{\ja{たたむ} (tatamu)}{たたむ (tatamu)}]{\ja{たたむ} (tatamu) — to fold}

\label{lesson:JA-C89-tatamu}



\begin{quote}
Before the new one: say the Japanese for to turn round, then the Japanese for to grab.
\end{quote}

\subsection*{You'll want to know: \ja{たたむ}}
\textbf{\ja{たたむ}} — \emph{tatamu} — \textquotedblleft{}to fold\textquotedblright{}.

\begin{grammarlens}[title={What you've built}]
One more word for this chapter.
\end{grammarlens}

\subsection*{Your turn}
\begin{itemize}
\raggedright
  \item \emph{Say it:} \emph{tatamu}
  \item \emph{Say it:} \emph{tatamu}, once more
  \item \emph{Say it:} say \emph{tatamu}
  \item \emph{Recall:} say the Japanese for to turn round, then the Japanese for to grab, then say \emph{tatamu} again
\end{itemize}

\begin{tcolorbox}[breakable,colback=teal!4,colframe=teal!35!black,title={Before you move on}]
What does \ja{たたむ} mean? (\textquotedblleft{}To fold\textquotedblright{}.) Say it once more.
\end{tcolorbox}
